% This template was tested with Pandoc 3.4 and pandoc-crossref v0.3.18.0. It should be backwards compatible with older version of pandoc..
\documentclass{article}


% if you need to pass options to natbib, use, e.g.:
%     \PassOptionsToPackage{numbers, compress}{natbib}
% before loading neurips_2023


% ready for submission
\usepackage[final,nonatbib]{neurips}


% to compile a preprint version, e.g., for submission to arXiv, add add the
% [preprint] option:
%     \usepackage[preprint]{neurips_2023}


% to compile a camera-ready version, add the [final] option, e.g.:
%     \usepackage[final]{neurips_2023}


% to avoid loading the natbib package, add option nonatbib:
%    \usepackage[nonatbib]{neurips_2023}


\usepackage[utf8]{inputenc} % allow utf-8 input
\usepackage[T1]{fontenc}    % use 8-bit T1 fonts
\usepackage{hyperref}       % hyperlinks
\usepackage{url}            % simple URL typesetting
\usepackage{booktabs}       % professional-quality tables
\usepackage{amsfonts}       % blackboard math symbols
\usepackage{nicefrac}       % compact symbols for 1/2, etc.
\usepackage{microtype}      % microtypography
\usepackage{xcolor}         % colors
\usepackage{graphicx}
\usepackage{longtable} % Add support for Pandoc's longtable if needed
\usepackage{array}     % For table alignment improvements
\usepackage{amsmath}
\usepackage{textcomp}
\setlength{\LTcapwidth}{\textwidth} % To make captions fit within page width

\makeatletter
\newsavebox\pandoc@box
\newcommand*\pandocbounded[1]{% scales image to fit in text height/width
  \sbox\pandoc@box{#1}%
  \Gscale@div\@tempa{\textheight}{\dimexpr\ht\pandoc@box+\dp\pandoc@box\relax}%
  \Gscale@div\@tempb{\linewidth}{\wd\pandoc@box}%
  \ifdim\@tempb\p@<\@tempa\p@\let\@tempa\@tempb\fi% select the smaller of both
  \ifdim\@tempa\p@<\p@\scalebox{\@tempa}{\usebox\pandoc@box}%
  \else\usebox{\pandoc@box}%
  \fi%
}
\makeatother
\makeatletter
\def\maxwidth{\ifdim\Gin@nat@width>\linewidth\linewidth\else\Gin@nat@width\fi}
\def\maxheight{\ifdim\Gin@nat@height>\textheight\textheight\else\Gin@nat@height\fi}
\makeatother
% Scale images if necessary, so that they will not overflow the page
% margins by default, and it is still possible to overwrite the defaults
% using explicit options in \includegraphics[width, height, ...]{}
\setkeys{Gin}{width=\maxwidth,height=\maxheight,keepaspectratio}
% Set default figure placement to htbp
\makeatletter
\def\fps@figure{htbp}
\makeatother

\providecommand{\tightlist}{%
  \setlength{\itemsep}{0pt}\setlength{\parskip}{0pt}}
\title{}


% Iterate through the authors except last to add \And. 

\author{%
}

% \author{%
%   David S.~Hippocampus \\
%   Department of Computer Science\\
%   Cranberry-Lemon University\\
%   Pittsburgh, PA 15213 \\
%   \texttt{hippo@cs.cranberry-lemon.edu} \\
%   % examples of more authors
%   % \And
%   % Coauthor \\
%   % Affiliation \\
%   % Address \\
%   % \texttt{email} \\
%   % \AND
%   % Coauthor \\
%   % Affiliation \\
%   % Address \\
%   % \texttt{email} \\
%   % \And
%   % Coauthor \\
%   % Affiliation \\
%   % Address \\
%   % \texttt{email} \\
%   % \And
%   % Coauthor \\
%   % Affiliation \\
%   % Address \\
%   % \texttt{email} \\
% }


\begin{document}


\maketitle


\begin{abstract}
    Add your abstract at the beginning of your markdown file like this 
  \begin{verbatim}
  --- 
  title: "Your Title" 
  abstract: "your abstract here"
  authors:
  - name: Leonardo V. Castorina
    affiliation: School of Informatics
    institution: University of Edinburgh
    email: justanemail@domain.ext
    address: Edinburgh
  - name: Coauthor
    affiliation: Affiliation
    institution: Institution
    email: coauthor@example.com
    address: Address
  ---
  \end{verbatim}
  This is called YAML frontmatter. If you set your abstract correctly you should not see this message.
  \end{abstract}


\section{Dimostrazioni di Machine Learning --- versione riorganizzata e
completa}\label{dimostrazioni-di-machine-learning-versione-riorganizzata-e-completa}

Questa dispensa riorganizza e migliora gli appunti forniti, con quattro
obiettivi:

\begin{enumerate}
\def\labelenumi{\arabic{enumi}.}
\tightlist
\item
  \textbf{notazione coerente} in tutto il testo;
\item
  \textbf{dimostrazioni passo passo}, senza salti logici;
\item
  \textbf{intuizione geometrica e statistica} accanto ai calcoli;
\item
  \textbf{correzione dei punti più facili da confondere} all'esame.
\end{enumerate}

\begin{center}\rule{0.5\linewidth}{0.5pt}\end{center}

\section{Convenzioni e notazione}\label{convenzioni-e-notazione}

Userò quasi sempre queste convenzioni:

\begin{itemize}
\tightlist
\item
  \(n\): numero di campioni
\item
  \(d\): numero di feature
\item
  \(x_i \in \mathbb{R}^d\): i-esimo input
\item
  \(y_i\): i-esima etichetta / target
\item
  \(X \in \mathbb{R}^{n \times d}\): \textbf{design matrix}
\item
  \(\theta \in \mathbb{R}^d\): vettore dei parametri
\item
  \(\hat y = X\theta\): vettore delle predizioni
\item
  \(r = y - X\theta\): vettore dei residui
\end{itemize}

Quando serve l'intercetta, la assorbo nella matrice \(X\) aggiungendo
una colonna di \(1\).

\begin{center}\rule{0.5\linewidth}{0.5pt}\end{center}

\section{1. Predittore ottimale con diverse
loss}\label{predittore-ottimale-con-diverse-loss}

L'idea generale è sempre la stessa: fissato un valore \(x\), vogliamo
scegliere una predizione \(z\) che minimizzi il \textbf{rischio
condizionato}

\[
\phi(z) = \mathbb{E}[\ell(z,Y)\mid X=x].
\]

Il predittore ottimale è quindi

\[
f^*(x) = \arg\min_z \mathbb{E}[\ell(z,Y)\mid X=x].
\]

\begin{center}\rule{0.5\linewidth}{0.5pt}\end{center}

\subsection{1.1 Perdita quadratica: il predittore ottimale è la media
condizionata}\label{perdita-quadratica-il-predittore-ottimale-uxe8-la-media-condizionata}

\subsubsection{Obiettivo}\label{obiettivo}

Dimostrare che, se

\[
\ell(z,Y)=(z-Y)^2,
\]

allora il predittore ottimale è

\[
f^*(x)=\mathbb{E}[Y\mid X=x].
\]

\subsubsection{Dimostrazione passo
passo}\label{dimostrazione-passo-passo}

Fissiamo \(x\) e poniamo

\[
\mu(x)=\mathbb{E}[Y\mid X=x].
\]

Dobbiamo minimizzare

\[
\phi(z)=\mathbb{E}[(z-Y)^2\mid X=x].
\]

Aggiungiamo e sottraiamo \(\mu(x)\) dentro la parentesi:

\[
z-Y = (z-\mu) + (\mu-Y).
\]

Quindi

\[
(z-Y)^2 = \big((z-\mu)+(\mu-Y)\big)^2.
\]

Sviluppiamo il quadrato:

\[
(z-Y)^2 = (z-\mu)^2 + (\mu-Y)^2 + 2(z-\mu)(\mu-Y).
\]

Ora facciamo il valore atteso condizionato:

\[
\phi(z)=\mathbb{E}[(z-\mu)^2\mid X=x]
+ \mathbb{E}[(\mu-Y)^2\mid X=x]
+ 2\mathbb{E}[(z-\mu)(\mu-Y)\mid X=x].
\]

Osserviamo che \(z-\mu\) è una costante rispetto a \(Y\) quando \(x\) è
fissato. Quindi:

\[
\phi(z)=(z-\mu)^2 + \mathbb{E}[(\mu-Y)^2\mid X=x]
+2(z-\mu)\mathbb{E}[\mu-Y\mid X=x].
\]

Ma

\[
\mathbb{E}[\mu-Y\mid X=x]
= \mu - \mathbb{E}[Y\mid X=x]
= \mu-\mu=0.
\]

Allora il termine misto sparisce e resta

\[
\phi(z)=(z-\mu)^2 + \mathbb{E}[(\mu-Y)^2\mid X=x].
\]

Il secondo termine non dipende da \(z\). Per minimizzare \(\phi(z)\),
dobbiamo quindi minimizzare solo \((z-\mu)^2\), che è minimo in
\(z=\mu\).

\subsubsection{Conclusione}\label{conclusione}

\[
f^*(x)=\mathbb{E}[Y\mid X=x].
\]

\subsubsection{Interpretazione}\label{interpretazione}

Con la loss quadratica, il miglior numero da predire è il
\textbf{baricentro} della distribuzione condizionata di \(Y\).

\begin{center}\rule{0.5\linewidth}{0.5pt}\end{center}

\subsection{1.2 Perdita assoluta: il predittore ottimale è la mediana
condizionata}\label{perdita-assoluta-il-predittore-ottimale-uxe8-la-mediana-condizionata}

\subsubsection{Obiettivo}\label{obiettivo-1}

Dimostrare che, se

\[
\ell(z,Y)=|z-Y|,
\]

allora il predittore ottimale è una \textbf{mediana condizionata} di
\(Y\mid X=x\).

\subsubsection{Dimostrazione passo
passo}\label{dimostrazione-passo-passo-1}

Fissiamo \(x\) e definiamo

\[
\phi(z)=\mathbb{E}[|z-Y|\mid X=x].
\]

Supponiamo per semplicità che \(Y\mid X=x\) abbia densità
\(p(y\mid x)\). Allora

\[
\phi(z)=\int_{-\infty}^{+\infty}|z-y|\,p(y\mid x)\,dy.
\]

Poiché il valore assoluto cambia forma in \(y=z\), separiamo
l'integrale:

\[
\phi(z)=\int_{-\infty}^{z}(z-y)p(y\mid x)\,dy
+\int_{z}^{+\infty}(y-z)p(y\mid x)\,dy.
\]

Deriviamo rispetto a \(z\) usando la regola di Leibniz.

\subsubsection{Derivata del primo
termine}\label{derivata-del-primo-termine}

\[
\frac{d}{dz}\int_{-\infty}^{z}(z-y)p(y\mid x)\,dy
=
\int_{-\infty}^{z}1\cdot p(y\mid x)\,dy + (z-z)p(z\mid x).
\]

Il termine al bordo è nullo, quindi resta

\[
\int_{-\infty}^{z}p(y\mid x)\,dy = P(Y\le z\mid X=x).
\]

\subsubsection{Derivata del secondo
termine}\label{derivata-del-secondo-termine}

\[
\frac{d}{dz}\int_{z}^{+\infty}(y-z)p(y\mid x)\,dy
=
\int_{z}^{+\infty}(-1)\,p(y\mid x)\,dy -(z-z)p(z\mid x).
\]

Anche qui il termine al bordo è nullo, quindi resta

\[
-\int_{z}^{+\infty}p(y\mid x)\,dy = -P(Y\ge z\mid X=x).
\]

\subsubsection{Condizione di
ottimalità}\label{condizione-di-ottimalituxe0}

Mettiendo insieme i due pezzi,

\[
\phi'(z)=P(Y\le z\mid X=x)-P(Y\ge z\mid X=x).
\]

Nel punto di minimo vogliamo \(\phi'(z)=0\), quindi

\[
P(Y\le z\mid X=x)=P(Y\ge z\mid X=x).
\]

Questa è precisamente la caratterizzazione della \textbf{mediana}.

\subsubsection{Conclusione}\label{conclusione-1}

\[
f^*(x)=\text{mediana di }Y\mid X=x.
\]

\subsubsection{Interpretazione}\label{interpretazione-1}

Con la loss assoluta non conta il ``centro di massa'' della
distribuzione, ma il punto che lascia metà probabilità a sinistra e metà
a destra.

\begin{center}\rule{0.5\linewidth}{0.5pt}\end{center}

\subsection{1.3 Classificazione binaria: classificatore di
Bayes}\label{classificazione-binaria-classificatore-di-bayes}

\subsubsection{Obiettivo}\label{obiettivo-2}

Supponiamo \(Y\in\{0,1\}\). Vogliamo capire quale decisione minimizza il
rischio condizionato.

\subsubsection{Caso generale: loss con costi
arbitrari}\label{caso-generale-loss-con-costi-arbitrari}

Se scegliamo l'azione \(a\in\{0,1\}\), il rischio condizionato è

\[
R(a\mid x)=\sum_{y\in\{0,1\}}\ell(a,y)P(Y=y\mid X=x).
\]

Quindi:

\begin{itemize}
\tightlist
\item
  se prediciamo \(1\),
\end{itemize}

\[
R(1\mid x)=\ell(1,0)P(Y=0\mid x)+\ell(1,1)P(Y=1\mid x);
\]

\begin{itemize}
\tightlist
\item
  se prediciamo \(0\),
\end{itemize}

\[
R(0\mid x)=\ell(0,0)P(Y=0\mid x)+\ell(0,1)P(Y=1\mid x).
\]

La regola ottimale è: predici \(1\) se \(R(1\mid x)\le R(0\mid x)\).

Sviluppiamo:

\[
\ell(1,0)P(Y=0\mid x)+\ell(1,1)P(Y=1\mid x)
\le
\ell(0,0)P(Y=0\mid x)+\ell(0,1)P(Y=1\mid x).
\]

Riordiniamo i termini:

\[
P(Y=1\mid x)\big(\ell(0,1)-\ell(1,1)\big)
\ge
P(Y=0\mid x)\big(\ell(1,0)-\ell(0,0)\big).
\]

Se dividiamo per \(P(Y=0\mid x)\) e per il denominatore positivo,
otteniamo una regola di soglia:

\[
\frac{P(Y=1\mid x)}{P(Y=0\mid x)}
\ge
\frac{\ell(1,0)-\ell(0,0)}{\ell(0,1)-\ell(1,1)}.
\]

\subsubsection{Caso particolare: 0-1
loss}\label{caso-particolare-0-1-loss}

Se la loss è quella standard di classificazione,

\[
\ell(\hat y,y)=\mathbf{1}\{\hat y\neq y\},
\]

allora:

\begin{itemize}
\tightlist
\item
  costo corretto \(=0\)
\item
  costo errore \(=1\)
\end{itemize}

Quindi:

\[
R(1\mid x)=P(Y=0\mid x), \qquad R(0\mid x)=P(Y=1\mid x).
\]

Prediciamo \(1\) se

\[
P(Y=0\mid x)\le P(Y=1\mid x),
\]

cioè

\[
P(Y=1\mid x)\ge \frac{1}{2}.
\]

\subsubsection{Conclusione}\label{conclusione-2}

Il classificatore ottimale è il \textbf{classificatore di Bayes}:

\[
f^*(x)=
\begin{cases}
1 & \text{se } P(Y=1\mid X=x)\ge 1/2,\\
0 & \text{altrimenti.}
\end{cases}
\]

\subsubsection{Interpretazione}\label{interpretazione-2}

Con loss 0-1, la decisione ottimale è ``scegli la classe più
probabile''.

\begin{center}\rule{0.5\linewidth}{0.5pt}\end{center}

\section{\texorpdfstring{2. Regressione lineare semplice: trovare \(m\)
e
\(q\)}{2. Regressione lineare semplice: trovare m e q}}\label{regressione-lineare-semplice-trovare-m-e-q}

Studiamo il modello

\[
\hat y_i = mx_i + q.
\]

Vogliamo minimizzare la loss quadratica empirica

\[
\hat R(m,q)=\frac{1}{n}\sum_{i=1}^n (y_i-mx_i-q)^2.
\]

Poiché la funzione è derivabile, cerchiamo il minimo imponendo che il
gradiente sia nullo.

\begin{center}\rule{0.5\linewidth}{0.5pt}\end{center}

\subsection{\texorpdfstring{2.1 Derivata rispetto a
\(q\)}{2.1 Derivata rispetto a q}}\label{derivata-rispetto-a-q}

Calcoliamo

\[
\frac{\partial \hat R}{\partial q}
=
\frac{1}{n}\sum_{i=1}^n 2(y_i-mx_i-q)(-1).
\]

Porre la derivata uguale a zero equivale a

\[
\sum_{i=1}^n (y_i-mx_i-q)=0.
\]

Separiamo i termini:

\[
\sum y_i - m\sum x_i - \sum q = 0.
\]

Dato che \(\sum q = nq\),

\[
\sum y_i - m\sum x_i - nq = 0.
\]

Dividendo per \(n\),

\[
\bar y - m\bar x - q = 0,
\]

quindi

\[
\hat q = \bar y - m\bar x.
\]

\subsubsection{Interpretazione}\label{interpretazione-3}

La retta ottima passa sempre per il punto medio del dataset
\((\bar x,\bar y)\).

\begin{center}\rule{0.5\linewidth}{0.5pt}\end{center}

\subsection{\texorpdfstring{2.2 Derivata rispetto a
\(m\)}{2.2 Derivata rispetto a m}}\label{derivata-rispetto-a-m}

Calcoliamo

\[
\frac{\partial \hat R}{\partial m}
=
\frac{1}{n}\sum_{i=1}^n 2(y_i-mx_i-q)(-x_i).
\]

Poniamo la derivata uguale a zero:

\[
\sum_{i=1}^n x_i(y_i-mx_i-q)=0.
\]

Sostituiamo \(q=\bar y - m\bar x\):

\[
\sum x_i\big(y_i-mx_i-(\bar y-m\bar x)\big)=0.
\]

Sviluppiamo:

\[
\sum x_i\big((y_i-\bar y)-m(x_i-\bar x)\big)=0.
\]

Portiamo fuori \(m\):

\[
\sum x_i(y_i-\bar y)-m\sum x_i(x_i-\bar x)=0.
\]

Da cui

\[
m=\frac{\sum x_i(y_i-\bar y)}{\sum x_i(x_i-\bar x)}.
\]

Questa formula è corretta, ma si può semplificare in una forma più
standard.

\begin{center}\rule{0.5\linewidth}{0.5pt}\end{center}

\subsection{2.3 Forma standard della
pendenza}\label{forma-standard-della-pendenza}

Usiamo il fatto che

\[
\sum_{i=1}^n (y_i-\bar y)=0.
\]

Allora

\[
\sum x_i(y_i-\bar y)
=
\sum (x_i-\bar x)(y_i-\bar y)
+
\bar x \sum (y_i-\bar y)
=
\sum (x_i-\bar x)(y_i-\bar y).
\]

Analogamente,

\[
\sum x_i(x_i-\bar x)
=
\sum (x_i-\bar x)^2.
\]

Quindi la formula standard è

\[
\hat m
=
\frac{\sum_{i=1}^n (x_i-\bar x)(y_i-\bar y)}
{\sum_{i=1}^n (x_i-\bar x)^2}.
\]

E quindi

\[
\hat q = \bar y - \hat m \bar x.
\]

\subsubsection{Forma statistica}\label{forma-statistica}

\[
\hat m = \frac{\operatorname{Cov}(X,Y)}{\operatorname{Var}(X)}.
\]

\begin{center}\rule{0.5\linewidth}{0.5pt}\end{center}

\section{3. Regressione lineare in forma matriciale e
OLS}\label{regressione-lineare-in-forma-matriciale-e-ols}

Passiamo al modello lineare generale

\[
\hat y = X\theta,
\]

dove

\begin{itemize}
\tightlist
\item
  \(X \in \mathbb{R}^{n\times d}\) è la design matrix;
\item
  \(\theta\in\mathbb{R}^d\) è il vettore dei parametri;
\item
  \(y\in\mathbb{R}^n\) è il vettore dei target.
\end{itemize}

Vogliamo minimizzare

\[
\hat R(\theta)=\frac{1}{n}\|y-X\theta\|_2^2.
\]

\begin{center}\rule{0.5\linewidth}{0.5pt}\end{center}

\subsection{3.1 Espansione della loss}\label{espansione-della-loss}

Usiamo la definizione di norma quadratica:

\[
\|y-X\theta\|_2^2 = (y-X\theta)^T(y-X\theta).
\]

Sviluppiamo:

\[
(y-X\theta)^T(y-X\theta)
=
y^Ty - y^TX\theta - \theta^TX^Ty + \theta^TX^TX\theta.
\]

Osserviamo che \(y^TX\theta\) è uno scalare, quindi coincide con il suo
trasposto:

\[
y^TX\theta = \theta^TX^Ty.
\]

Allora

\[
\hat R(\theta)=\frac{1}{n}\big(y^Ty - 2\theta^TX^Ty + \theta^TX^TX\theta\big).
\]

\begin{center}\rule{0.5\linewidth}{0.5pt}\end{center}

\subsection{3.2 Gradiente e equazioni
normali}\label{gradiente-e-equazioni-normali}

Deriviamo rispetto a \(\theta\):

\begin{itemize}
\tightlist
\item
  \(\nabla_\theta(y^Ty)=0\), perché non dipende da \(\theta\);
\item
  \(\nabla_\theta(-2\theta^TX^Ty)=-2X^Ty\);
\item
  \(\nabla_\theta(\theta^TX^TX\theta)=2X^TX\theta\).
\end{itemize}

Quindi

\[
\nabla \hat R(\theta)=\frac{1}{n}\big(-2X^Ty+2X^TX\theta\big).
\]

Ponendo il gradiente uguale a zero:

\[
X^TX\theta=X^Ty.
\]

Queste sono le \textbf{equazioni normali}.

Se \(X^TX\) è invertibile, allora

\[
\hat\theta = (X^TX)^{-1}X^Ty.
\]

Questa è la soluzione \textbf{OLS (Ordinary Least Squares)}.

\begin{center}\rule{0.5\linewidth}{0.5pt}\end{center}

\subsection{\texorpdfstring{3.3 Perché compare
\(X^T\)?}{3.3 Perché compare X\^{}T?}}\label{perchuxe9-compare-xt}

Questa è una domanda d'esame molto frequente.

La design matrix \textbf{non è ``trasposta di suo''}. La trasposta
compare per due motivi:

\begin{enumerate}
\def\labelenumi{\arabic{enumi}.}
\tightlist
\item
  quando sviluppi il quadrato \((y-X\theta)^T(y-X\theta)\);
\item
  quando imponi la condizione geometrica di ortogonalità dei residui.
\end{enumerate}

Più precisamente, al minimo il residuo

\[
r = y-X\hat\theta
\]

deve essere ortogonale a tutte le colonne di \(X\). In forma matriciale:

\[
X^Tr=0.
\]

Sostituendo \(r\),

\[
X^T(y-X\hat\theta)=0
\quad\Longrightarrow\quad
X^TX\hat\theta=X^Ty.
\]

Quindi \(X^T\) compare perché ``trasporta'' l'informazione dallo spazio
delle osservazioni allo spazio delle feature, e perché codifica
l'ortogonalità del residuo alle colonne di \(X\).

\begin{center}\rule{0.5\linewidth}{0.5pt}\end{center}

\section{4. Interpretazione geometrica: proiezione
ortogonale}\label{interpretazione-geometrica-proiezione-ortogonale}

Una delle letture più importanti di OLS è geometrica.

Il vettore \(\hat y = X\hat\theta\) appartiene allo spazio generato
dalle colonne di \(X\), cioè

\[
\hat y \in \operatorname{span}(X).
\]

L'OLS sceglie il punto di \(\operatorname{span}(X)\) più vicino a \(y\)
in norma euclidea.

In altre parole, \(\hat y\) è la \textbf{proiezione ortogonale} di \(y\)
sul sottospazio \(\operatorname{span}(X)\).

\begin{center}\rule{0.5\linewidth}{0.5pt}\end{center}

\subsection{4.1 Matrice di proiezione}\label{matrice-di-proiezione}

Definiamo

\[
P = X(X^TX)^{-1}X^T.
\]

Allora

\[
\hat y = Py.
\]

Vogliamo dimostrare che \(P\) è una proiezione ortogonale.

\begin{center}\rule{0.5\linewidth}{0.5pt}\end{center}

\subsection{4.2 Prima proprietà: ogni vettore nello span resta
invariato}\label{prima-proprietuxe0-ogni-vettore-nello-span-resta-invariato}

Sia \(v\in\operatorname{span}(X)\). Allora esiste \(\lambda\) tale che

\[
v=X\lambda.
\]

Applichiamo \(P\):

\[
Pv = X(X^TX)^{-1}X^T(X\lambda).
\]

Raggruppiamo i termini:

\[
Pv = X(X^TX)^{-1}(X^TX)\lambda = XI\lambda = X\lambda = v.
\]

Quindi \(P\) lascia invariati i vettori che stanno già nello spazio.

\begin{center}\rule{0.5\linewidth}{0.5pt}\end{center}

\subsection{4.3 Seconda proprietà: ogni vettore ortogonale allo span
viene mandato a
zero}\label{seconda-proprietuxe0-ogni-vettore-ortogonale-allo-span-viene-mandato-a-zero}

Sia \(u\) tale che \(u\perp \operatorname{span}(X)\). Questo equivale a
dire

\[
X^Tu=0.
\]

Allora

\[
Pu = X(X^TX)^{-1}X^Tu = X(X^TX)^{-1}0 = 0.
\]

Quindi \(P\) annulla i vettori ortogonali allo spazio.

\begin{center}\rule{0.5\linewidth}{0.5pt}\end{center}

\subsection{4.4 Conclusione geometrica}\label{conclusione-geometrica}

Le due proprietà precedenti caratterizzano una proiezione ortogonale.
Quindi

\[
\hat y = Py
\]

è la proiezione ortogonale di \(y\) su \(\operatorname{span}(X)\).

\subsubsection{Conseguenza chiave}\label{conseguenza-chiave}

Il residuo

\[
r=y-\hat y
\]

è ortogonale allo spazio generato dalle colonne di \(X\), cioè

\[
X^Tr=0.
\]

Questa è esattamente la stessa informazione contenuta nelle equazioni
normali.

\begin{center}\rule{0.5\linewidth}{0.5pt}\end{center}

\section{5. Ottimalità di OLS e
Hessiana}\label{ottimalituxe0-di-ols-e-hessiana}

Per capire perché il punto trovato è davvero un minimo globale,
guardiamo la derivata seconda.

Partiamo da

\[
\hat R(\theta)=\frac{1}{2n}\|X\theta-y\|_2^2.
\]

Uso il fattore \(1/2\) solo per semplificare le derivate.

\begin{center}\rule{0.5\linewidth}{0.5pt}\end{center}

\subsection{5.1 Gradiente}\label{gradiente}

\[
\nabla \hat R(\theta)=\frac{1}{n}X^T(X\theta-y).
\]

\begin{center}\rule{0.5\linewidth}{0.5pt}\end{center}

\subsection{5.2 Hessiana}\label{hessiana}

Derivando ancora rispetto a \(\theta\),

\[
H(\theta)=\nabla^2 \hat R(\theta)=\frac{1}{n}X^TX.
\]

Questa Hessiana non dipende nemmeno da \(\theta\): è costante.

\begin{center}\rule{0.5\linewidth}{0.5pt}\end{center}

\subsection{5.3 Perché è semidefinita
positiva}\label{perchuxe9-uxe8-semidefinita-positiva}

Per dimostrare che \(H\) è semidefinita positiva, prendiamo un vettore
arbitrario \(v\in\mathbb{R}^d\) e calcoliamo

\[
v^THv = \frac{1}{n}v^TX^TXv.
\]

Osserviamo che

\[
v^TX^T = (Xv)^T.
\]

Quindi

\[
v^THv = \frac{1}{n}(Xv)^T(Xv)=\frac{1}{n}\|Xv\|_2^2 \ge 0.
\]

Poiché è sempre non negativo, \(H\) è semidefinita positiva.

\subsubsection{Conclusione}\label{conclusione-3}

La funzione di costo è \textbf{convessa}. Quindi ogni minimo locale è
anche un \textbf{minimo globale}.

Se inoltre \(X^TX\) è definita positiva (per esempio se le colonne di
\(X\) sono linearmente indipendenti), il minimo è unico.

\begin{center}\rule{0.5\linewidth}{0.5pt}\end{center}

\section{6. Proprietà statistiche di
OLS}\label{proprietuxe0-statistiche-di-ols}

Assumiamo il modello lineare corretto

\[
y = X\theta^* + \varepsilon,
\]

dove

\begin{itemize}
\tightlist
\item
  \(\theta^*\) è il parametro vero;
\item
  \(\varepsilon\) è il rumore con \[
  \mathbb{E}[\varepsilon]=0,\qquad \mathbb{E}[\varepsilon\varepsilon^T]=\sigma^2 I.
  \]
\end{itemize}

\begin{center}\rule{0.5\linewidth}{0.5pt}\end{center}

\subsection{6.1 Correttezza (unbiasedness) di
OLS}\label{correttezza-unbiasedness-di-ols}

Vogliamo mostrare che

\[
\mathbb{E}[\hat\theta]=\theta^*.
\]

Partiamo dalla formula di OLS:

\[
\hat\theta=(X^TX)^{-1}X^Ty.
\]

Sostituiamo \(y=X\theta^*+\varepsilon\):

\[
\hat\theta=(X^TX)^{-1}X^T(X\theta^*+\varepsilon).
\]

Distribuiamo:

\[
\hat\theta=(X^TX)^{-1}X^TX\theta^* + (X^TX)^{-1}X^T\varepsilon.
\]

Il primo termine si semplifica:

\[
\hat\theta=\theta^* + (X^TX)^{-1}X^T\varepsilon.
\]

Ora facciamo il valore atteso:

\[
\mathbb{E}[\hat\theta]
=
\theta^* + (X^TX)^{-1}X^T\mathbb{E}[\varepsilon].
\]

Poiché \(\mathbb{E}[\varepsilon]=0\),

\[
\mathbb{E}[\hat\theta]=\theta^*.
\]

\subsubsection{Conclusione}\label{conclusione-4}

OLS è \textbf{unbiased}: in media non introduce errore sistematico sui
parametri.

\begin{center}\rule{0.5\linewidth}{0.5pt}\end{center}

\subsection{6.2 Varianza dello stimatore
OLS}\label{varianza-dello-stimatore-ols}

Vogliamo calcolare

\[
\operatorname{Var}(\hat\theta)
=
\mathbb{E}\big[(\hat\theta-\theta^*)(\hat\theta-\theta^*)^T\big].
\]

Dalla formula appena ricavata,

\[
\hat\theta-\theta^*=(X^TX)^{-1}X^T\varepsilon.
\]

Allora

\[
\operatorname{Var}(\hat\theta)
=
\mathbb{E}\Big[
\big((X^TX)^{-1}X^T\varepsilon\big)
\big((X^TX)^{-1}X^T\varepsilon\big)^T
\Big].
\]

Ora usiamo la regola del trasposto \((ABC)^T=C^TB^TA^T\):

\[
\big((X^TX)^{-1}X^T\varepsilon\big)^T
=
\varepsilon^TX\big((X^TX)^{-1}\big)^T.
\]

Poiché \(X^TX\) è simmetrica, anche la sua inversa è simmetrica:

\[
\big((X^TX)^{-1}\big)^T=(X^TX)^{-1}.
\]

Quindi

\[
\operatorname{Var}(\hat\theta)
=
\mathbb{E}\Big[
(X^TX)^{-1}X^T\varepsilon\varepsilon^TX(X^TX)^{-1}
\Big].
\]

Portiamo fuori dall'aspettativa tutto ciò che è deterministico:

\[
\operatorname{Var}(\hat\theta)
=
(X^TX)^{-1}X^T
\mathbb{E}[\varepsilon\varepsilon^T]
X(X^TX)^{-1}.
\]

Usiamo l'ipotesi \(\mathbb{E}[\varepsilon\varepsilon^T]=\sigma^2 I\):

\[
\operatorname{Var}(\hat\theta)
=
(X^TX)^{-1}X^T(\sigma^2 I)X(X^TX)^{-1}.
\]

Portiamo fuori \(\sigma^2\):

\[
\operatorname{Var}(\hat\theta)
=
\sigma^2(X^TX)^{-1}X^TX(X^TX)^{-1}.
\]

Semplificando,

\[
\operatorname{Var}(\hat\theta)=\sigma^2(X^TX)^{-1}.
\]

\subsubsection{Interpretazione}\label{interpretazione-4}

\begin{itemize}
\tightlist
\item
  più rumore (\(\sigma^2\) alto) \(\Rightarrow\) più incertezza;
\item
  più informazione nei dati (\(X^TX\) grande) \(\Rightarrow\) varianza
  più piccola.
\end{itemize}

\begin{center}\rule{0.5\linewidth}{0.5pt}\end{center}

\section{7. Scomposizione del rischio per la regressione
lineare}\label{scomposizione-del-rischio-per-la-regressione-lineare}

Supponiamo ancora che il modello vero sia

\[
Y = \phi(X)^T\theta^* + \varepsilon,
\]

con \(\mathbb{E}[\varepsilon\mid X]=0\) e
\(\operatorname{Var}(\varepsilon\mid X)=\sigma^2\).

Definiamo il rischio di un parametro \(\theta\) come

\[
R(\theta)=\mathbb{E}\big[(\phi(X)^T\theta - Y)^2\big].
\]

Vogliamo dimostrare

\[
R(\theta)=
(\theta-\theta^*)^T\Sigma(\theta-\theta^*)+\sigma^2,
\]

dove

\[
\Sigma=\mathbb{E}[\phi(X)\phi(X)^T].
\]

\begin{center}\rule{0.5\linewidth}{0.5pt}\end{center}

\subsection{7.1 Dimostrazione passo
passo}\label{dimostrazione-passo-passo-2}

Sostituiamo \(Y=\phi(X)^T\theta^*+\varepsilon\):

\[
R(\theta)
=
\mathbb{E}\Big[
\big(\phi(X)^T\theta - \phi(X)^T\theta^* - \varepsilon\big)^2
\Big].
\]

Raggruppiamo:

\[
R(\theta)
=
\mathbb{E}\Big[
\big(\phi(X)^T(\theta-\theta^*)-\varepsilon\big)^2
\Big].
\]

Sviluppiamo il quadrato:

\[
R(\theta)
=
\mathbb{E}\Big[
(\phi(X)^T(\theta-\theta^*))^2
+ \varepsilon^2
-2\phi(X)^T(\theta-\theta^*)\varepsilon
\Big].
\]

Separiamo i tre termini.

\subsubsection{Primo termine}\label{primo-termine}

\[
(\phi(X)^T(\theta-\theta^*))^2
=
(\theta-\theta^*)^T\phi(X)\phi(X)^T(\theta-\theta^*).
\]

Prendendo il valore atteso,

\[
\mathbb{E}\big[(\phi(X)^T(\theta-\theta^*))^2\big]
=
(\theta-\theta^*)^T
\mathbb{E}[\phi(X)\phi(X)^T]
(\theta-\theta^*)
=
(\theta-\theta^*)^T\Sigma(\theta-\theta^*).
\]

\subsubsection{Secondo termine}\label{secondo-termine}

\[
\mathbb{E}[\varepsilon^2]=\sigma^2.
\]

\subsubsection{Terzo termine}\label{terzo-termine}

Poiché \(\mathbb{E}[\varepsilon\mid X]=0\),

\[
\mathbb{E}\big[\phi(X)^T(\theta-\theta^*)\varepsilon\big]=0.
\]

Quindi il termine incrociato scompare.

\subsubsection{Conclusione}\label{conclusione-5}

\[
R(\theta)=
(\theta-\theta^*)^T\Sigma(\theta-\theta^*)+\sigma^2.
\]

\begin{center}\rule{0.5\linewidth}{0.5pt}\end{center}

\subsection{7.2 Interpretazione}\label{interpretazione-5}

Il rischio ha due pezzi:

\begin{enumerate}
\def\labelenumi{\arabic{enumi}.}
\item
  \textbf{errore di stima / approssimazione} \[
  (\theta-\theta^*)^T\Sigma(\theta-\theta^*)
  \] che dipende da quanto siamo lontani dal parametro ottimo;
\item
  \textbf{rumore irriducibile} \[
  \sigma^2,
  \] che non possiamo eliminare nemmeno con il modello perfetto.
\end{enumerate}

\begin{center}\rule{0.5\linewidth}{0.5pt}\end{center}

\section{8. Overfitting, bias-variance e Ridge
Regression}\label{overfitting-bias-variance-e-ridge-regression}

Quando il modello è troppo flessibile rispetto ai dati disponibili, può
adattarsi anche al rumore del training set: questo è
l'\textbf{overfitting}.

\begin{center}\rule{0.5\linewidth}{0.5pt}\end{center}

\subsection{8.1 Bias e varianza: idea
intuitiva}\label{bias-e-varianza-idea-intuitiva}

\begin{itemize}
\tightlist
\item
  \textbf{Bias alto}: il modello è troppo semplice e sbaglia in modo
  sistematico.
\item
  \textbf{Varianza alta}: il modello cambia molto se cambiano i dati di
  training.
\end{itemize}

Aumentando la complessità del modello, in generale:

\begin{itemize}
\tightlist
\item
  il bias tende a scendere;
\item
  la varianza tende a salire.
\end{itemize}

L'overfitting è proprio il regime in cui la varianza domina.

\begin{center}\rule{0.5\linewidth}{0.5pt}\end{center}

\subsection{8.2 Ridge Regression}\label{ridge-regression}

Per combattere l'overfitting, aggiungiamo una penalizzazione \(L_2\):

\[
J(\theta)=\frac{1}{n}\|y-X\theta\|_2^2+\lambda\|\theta\|_2^2,
\qquad \lambda\ge 0.
\]

\subsubsection{Calcolo del gradiente}\label{calcolo-del-gradiente}

\[
\nabla J(\theta)
=
\frac{2}{n}X^T(X\theta-y)+2\lambda\theta.
\]

Ponendo il gradiente uguale a zero,

\[
\frac{1}{n}X^TX\theta - \frac{1}{n}X^Ty + \lambda\theta = 0.
\]

Moltiplichiamo per \(n\):

\[
X^TX\theta - X^Ty + n\lambda\theta = 0.
\]

Equivalentemente, assorbendo il fattore \(n\) dentro \(\lambda\) se la
convenzione del corso lo fa, si ottiene la forma più comune

\[
(X^TX+\lambda I)\theta = X^Ty.
\]

Quindi

\[
\hat\theta_{\text{ridge}}=(X^TX+\lambda I)^{-1}X^Ty.
\]

\subsubsection{Perché aiuta}\label{perchuxe9-aiuta}

La matrice \(X^TX+\lambda I\) è più stabile numericamente e, per
\(\lambda>0\), risulta invertibile anche in situazioni problematiche in
cui \(X^TX\) è singolare o quasi singolare.

\subsubsection{Effetto statistico}\label{effetto-statistico}

Ridge introduce un po' di bias, ma riduce la varianza. Spesso il
bilanciamento migliora l'errore di generalizzazione.

\begin{center}\rule{0.5\linewidth}{0.5pt}\end{center}

\section{9. Complessità computazionale della soluzione
OLS}\label{complessituxe0-computazionale-della-soluzione-ols}

La soluzione chiusa è

\[
\hat\theta=(X^TX)^{-1}X^Ty.
\]

Supponiamo \(X\in\mathbb{R}^{n\times d}\).

\begin{center}\rule{0.5\linewidth}{0.5pt}\end{center}

\subsection{\texorpdfstring{9.1 Costo di
\(X^TX\)}{9.1 Costo di X\^{}TX}}\label{costo-di-xtx}

\(X^T\) è \(d\times n\), quindi \(X^TX\) è \(d\times d\).

Il costo della moltiplicazione è

\[
O(nd^2).
\]

\begin{center}\rule{0.5\linewidth}{0.5pt}\end{center}

\subsection{9.2 Costo dell'inversione}\label{costo-dellinversione}

Invertire una matrice \(d\times d\) costa, in prima approssimazione,

\[
O(d^3).
\]

\begin{center}\rule{0.5\linewidth}{0.5pt}\end{center}

\subsection{\texorpdfstring{9.3 Costo di
\(X^Ty\)}{9.3 Costo di X\^{}Ty}}\label{costo-di-xty}

La moltiplicazione matrice-vettore costa

\[
O(nd).
\]

\begin{center}\rule{0.5\linewidth}{0.5pt}\end{center}

\subsection{9.4 Complessità totale}\label{complessituxe0-totale}

\[
O(nd^2 + d^3).
\]

\subsubsection{Conseguenza pratica}\label{conseguenza-pratica}

Se \(d\) è molto grande, il termine \(d^3\) può diventare proibitivo.
Per questo si usano metodi iterativi come il \textbf{Gradient Descent}.

\begin{center}\rule{0.5\linewidth}{0.5pt}\end{center}

\section{10. Gradient Descent per la regressione
lineare}\label{gradient-descent-per-la-regressione-lineare}

Partiamo dalla loss

\[
\hat R(\theta)=\frac{1}{2n}\|X\theta-y\|_2^2.
\]

\begin{center}\rule{0.5\linewidth}{0.5pt}\end{center}

\subsection{10.1 Gradiente}\label{gradiente-1}

Abbiamo già visto che

\[
\nabla \hat R(\theta)=\frac{1}{n}X^T(X\theta-y).
\]

\begin{center}\rule{0.5\linewidth}{0.5pt}\end{center}

\subsection{10.2 Aggiornamento}\label{aggiornamento}

La regola del Gradient Descent è

\[
\theta_{t+1}=\theta_t-\gamma_t \nabla \hat R(\theta_t),
\]

quindi

\[
\theta_{t+1}
=
\theta_t - \frac{\gamma_t}{n}X^T(X\theta_t-y).
\]

Qui \(\gamma_t\) è il \textbf{learning rate}.

\begin{center}\rule{0.5\linewidth}{0.5pt}\end{center}

\subsection{10.3 Perché converge al minimo
globale}\label{perchuxe9-converge-al-minimo-globale}

Lo abbiamo già visto con l'Hessiana:

\[
\nabla^2 \hat R(\theta)=\frac{1}{n}X^TX \succeq 0.
\]

La loss è convessa, quindi non ci sono minimi locali sbagliati. Con uno
step size adeguato, Gradient Descent converge al minimo globale.

\begin{center}\rule{0.5\linewidth}{0.5pt}\end{center}

\subsection{10.4 Interpretazione dell'Hessiana e condition
number}\label{interpretazione-dellhessiana-e-condition-number}

Se \(H=\frac{1}{n}X^TX\), gli autovalori di \(H\) controllano la
curvatura della loss.

\begin{itemize}
\tightlist
\item
  Se gli autovalori sono ben bilanciati, la valle è ``rotonda'' e la
  convergenza è veloce.
\item
  Se il rapporto tra autovalore massimo e minimo è grande, la valle è
  stretta e allungata: il Gradient Descent zigzaga e converge
  lentamente.
\end{itemize}

Questo rapporto è il \textbf{condition number}

\[
\kappa = \frac{L}{\mu},
\]

dove \(L\) è il massimo autovalore e \(\mu\) il minimo positivo.

\begin{center}\rule{0.5\linewidth}{0.5pt}\end{center}

\section{11. Stochastic Gradient Descent
(SGD)}\label{stochastic-gradient-descent-sgd}

Il Gradient Descent batch usa tutto il dataset ad ogni iterazione. L'SGD
usa invece un solo esempio casuale o un mini-batch.

\begin{center}\rule{0.5\linewidth}{0.5pt}\end{center}

\subsection{11.1 Gradiente vero}\label{gradiente-vero}

Il gradiente della loss empirica può essere scritto come media dei
gradienti dei singoli punti:

\[
\nabla \hat R(\theta)
=
\frac{1}{n}\sum_{i=1}^n (\phi(x_i)^T\theta - y_i)\phi(x_i).
\]

\begin{center}\rule{0.5\linewidth}{0.5pt}\end{center}

\subsection{11.2 Stimatore stocastico}\label{stimatore-stocastico}

Scegliamo un indice \(i\) uniformemente a caso e definiamo

\[
\tilde g(\theta)
=
(\phi(x_i)^T\theta-y_i)\phi(x_i).
\]

Questo è un gradiente rumoroso, ma molto economico da calcolare.

\begin{center}\rule{0.5\linewidth}{0.5pt}\end{center}

\subsection{11.3 Perché è unbiased}\label{perchuxe9-uxe8-unbiased}

Calcoliamo il valore atteso rispetto alla scelta casuale dell'indice:

\[
\mathbb{E}[\tilde g(\theta)]
=
\sum_{j=1}^n P(i=j)\,(\phi(x_j)^T\theta-y_j)\phi(x_j).
\]

Poiché la scelta è uniforme, \(P(i=j)=1/n\). Quindi

\[
\mathbb{E}[\tilde g(\theta)]
=
\frac{1}{n}\sum_{j=1}^n (\phi(x_j)^T\theta-y_j)\phi(x_j)
=
\nabla \hat R(\theta).
\]

\subsubsection{Conclusione}\label{conclusione-6}

L'SGD è uno stimatore \textbf{corretto (unbiased)} del gradiente vero:
in media punta nella direzione giusta.

\begin{center}\rule{0.5\linewidth}{0.5pt}\end{center}

\section{12. Logistic Regression e Maximum Likelihood
Estimation}\label{logistic-regression-e-maximum-likelihood-estimation}

La Logistic Regression non nasce da una loss quadratica, ma da un
\textbf{modello probabilistico}.

Per \(Y\in\{0,1\}\), assumiamo

\[
P(Y=1\mid x)=\sigma(\theta^T\phi(x))
=
\frac{1}{1+e^{-\theta^T\phi(x)}}.
\]

Chiamiamo

\[
p_i = P(Y=1\mid x_i)=\sigma(\theta^T\phi(x_i)).
\]

Allora

\[
P(Y=0\mid x_i)=1-p_i.
\]

\begin{center}\rule{0.5\linewidth}{0.5pt}\end{center}

\subsection{12.1 Perché è collegata alla
MLE}\label{perchuxe9-uxe8-collegata-alla-mle}

Dato \(x_i\), la variabile \(Y_i\) viene modellata come una Bernoulli di
parametro \(p_i\):

\[
Y_i\mid x_i \sim \text{Bernoulli}(p_i).
\]

Per una variabile Bernoulli, la probabilità del valore osservato
\(y_i\in\{0,1\}\) si scrive in forma compatta come

\[
P(Y_i=y_i\mid x_i)=p_i^{y_i}(1-p_i)^{1-y_i}.
\]

Assumendo indipendenza dei campioni, la likelihood dell'intero dataset è

\[
L(\theta)=\prod_{i=1}^n p_i^{y_i}(1-p_i)^{1-y_i}.
\]

La \textbf{Maximum Likelihood Estimation} sceglie \(\theta\) che
massimizza \(L(\theta)\).

\begin{center}\rule{0.5\linewidth}{0.5pt}\end{center}

\subsection{12.2 Perché si usa il
log-likelihood}\label{perchuxe9-si-usa-il-log-likelihood}

Massimizzare un prodotto è scomodo. Applichiamo il logaritmo:

\[
\ell(\theta)=\log L(\theta)
=
\sum_{i=1}^n \Big(y_i\log p_i + (1-y_i)\log(1-p_i)\Big).
\]

Poiché il logaritmo è strettamente crescente, massimizzare \(L\) o
\(\ell\) è equivalente.

\begin{center}\rule{0.5\linewidth}{0.5pt}\end{center}

\subsection{12.3 Perché sembrano esistere due obiettivi
diversi}\label{perchuxe9-sembrano-esistere-due-obiettivi-diversi}

In pratica si vedono due formulazioni:

\subsubsection{Formulazione statistica}\label{formulazione-statistica}

\[
\max_\theta \ell(\theta).
\]

\subsubsection{Formulazione di ottimizzazione / machine
learning}\label{formulazione-di-ottimizzazione-machine-learning}

\[
\min_\theta -\ell(\theta).
\]

La quantità

\[
-\ell(\theta)
=
-\sum_{i=1}^n \Big(y_i\log p_i + (1-y_i)\log(1-p_i)\Big)
\]

si chiama:

\begin{itemize}
\tightlist
\item
  \textbf{negative log-likelihood}
\item
  \textbf{log-loss}
\item
  \textbf{binary cross-entropy}
\end{itemize}

Sono tre nomi per lo stesso oggetto.

\subsubsection{Messaggio da ricordare}\label{messaggio-da-ricordare}

Non ci sono due problemi diversi: è lo \textbf{stesso problema}, scritto
una volta come massimizzazione e una volta come minimizzazione.

\begin{center}\rule{0.5\linewidth}{0.5pt}\end{center}

\subsection{12.4 Perché non si ottimizza direttamente
l'accuracy}\label{perchuxe9-non-si-ottimizza-direttamente-laccuracy}

L'accuracy dipende da una soglia e non è liscia. Piccole variazioni di
probabilità possono non cambiare nulla o cambiare tutto di colpo.

La log-loss invece:

\begin{itemize}
\tightlist
\item
  è continua;
\item
  è derivabile;
\item
  penalizza molto le predizioni sbagliate ma molto sicure.
\end{itemize}

Per questo è adatta all'ottimizzazione con gradient methods.

\begin{center}\rule{0.5\linewidth}{0.5pt}\end{center}

\section{13. Gradiente e concavità della log-likelihood
logistica}\label{gradiente-e-concavituxe0-della-log-likelihood-logistica}

\begin{center}\rule{0.5\linewidth}{0.5pt}\end{center}

\subsection{13.1 Derivata della sigmoide}\label{derivata-della-sigmoide}

La sigmoide è

\[
\sigma(z)=\frac{1}{1+e^{-z}}.
\]

Deriviamo:

\[
\sigma'(z)
=
\frac{e^{-z}}{(1+e^{-z})^2}.
\]

Riscriviamo il risultato in forma elegante. Notiamo che

\[
\sigma(z)=\frac{1}{1+e^{-z}},
\qquad
1-\sigma(z)=\frac{e^{-z}}{1+e^{-z}}.
\]

Moltiplicandole,

\[
\sigma(z)(1-\sigma(z))
=
\frac{1}{1+e^{-z}}
\cdot
\frac{e^{-z}}{1+e^{-z}}
=
\frac{e^{-z}}{(1+e^{-z})^2}
=
\sigma'(z).
\]

Quindi la formula da ricordare è

\[
\sigma'(z)=\sigma(z)(1-\sigma(z)).
\]

\begin{center}\rule{0.5\linewidth}{0.5pt}\end{center}

\subsection{13.2 Gradiente del
log-likelihood}\label{gradiente-del-log-likelihood}

Per comodità poniamo \(h_\theta(x_i)=p_i=\sigma(\theta^T\phi(x_i))\).

Il log-likelihood è

\[
\ell(\theta)=\sum_{i=1}^n\Big(y_i\log p_i + (1-y_i)\log(1-p_i)\Big).
\]

Derivando rispetto a \(\theta\), si ottiene

\[
\nabla \ell(\theta)
=
\sum_{i=1}^n (y_i-p_i)\phi(x_i).
\]

In forma matriciale:

\[
\nabla \ell(\theta)=\Phi^T(y-p),
\]

dove \(p=(p_1,\dots,p_n)^T\).

Se invece minimizziamo la negative log-likelihood, il gradiente cambia
segno:

\[
\nabla(-\ell(\theta))=\Phi^T(p-y).
\]

\begin{center}\rule{0.5\linewidth}{0.5pt}\end{center}

\subsection{13.3 Hessiana e concavità}\label{hessiana-e-concavituxe0}

Derivando ancora, si ottiene

\[
H(\theta)=\nabla^2 \ell(\theta)
=
-\Phi^T W \Phi,
\]

dove \(W\) è la matrice diagonale

\[
W = \operatorname{diag}(p_i(1-p_i)).
\]

Poiché \(0\le p_i(1-p_i)\), la matrice \(W\) è semidefinita positiva.

Prendiamo un vettore arbitrario \(v\). Allora

\[
v^THv
=
-v^T\Phi^T W \Phi v
=
-(\Phi v)^T W (\Phi v).
\]

Chiamiamo \(u=\Phi v\). Allora

\[
v^THv = -u^TWu = -\sum_{i=1}^n w_i u_i^2 \le 0.
\]

Quindi l'Hessiana è semidefinita \textbf{negativa}.

\subsubsection{Conclusione}\label{conclusione-7}

Il log-likelihood \(\ell(\theta)\) è \textbf{concavo}. Dunque ha un
unico massimo globale (a meno di degenerazioni).

Se invece consideriamo la negative log-likelihood, la Hessiana cambia
segno e diventa semidefinita positiva: la loss è convessa.

\begin{center}\rule{0.5\linewidth}{0.5pt}\end{center}

\section{14. Legame tra cross-entropy e KL
divergence}\label{legame-tra-cross-entropy-e-kl-divergence}

Sia \(p\) la distribuzione vera e \(q_\theta\) quella predetta dal
modello.

La KL divergence è

\[
\mathrm{KL}(p\|q_\theta)=\sum_i p_i \log\frac{p_i}{q_{\theta,i}}.
\]

Sviluppiamo:

\[
\mathrm{KL}(p\|q_\theta)
=
\sum_i p_i\log p_i - \sum_i p_i\log q_{\theta,i}.
\]

Il primo termine non dipende da \(\theta\). Quindi minimizzare la KL
rispetto a \(\theta\) equivale a minimizzare

\[
-\sum_i p_i\log q_{\theta,i},
\]

che è precisamente la \textbf{cross-entropy}.

Perciò, in classificazione probabilistica:

\begin{itemize}
\tightlist
\item
  massimizzare la log-likelihood;
\item
  minimizzare la negative log-likelihood;
\item
  minimizzare la cross-entropy;
\item
  minimizzare la KL dalla distribuzione vera alla distribuzione del
  modello
\end{itemize}

sono tutte formulazioni equivalenti, a meno di costanti additive.

\begin{center}\rule{0.5\linewidth}{0.5pt}\end{center}

\section{15. Softmax Regression:
gradiente}\label{softmax-regression-gradiente}

Per il caso multiclasse con \(K\) classi, definiamo per la classe \(k\):

\[
s_k(x)=\phi(x)^T\theta^{(k)}.
\]

La Softmax definisce

\[
\hat P^{(k)}(x)
=
\frac{e^{s_k(x)}}{\sum_{j=1}^K e^{s_j(x)}}.
\]

Usiamo etichette one-hot \(y_i^{(k)}\in\{0,1\}\), con
\(\sum_k y_i^{(k)}=1\).

Il log-likelihood è

\[
L(\theta)
=
\sum_{i=1}^n\sum_{k=1}^K y_i^{(k)}\log \hat P^{(k)}(x_i).
\]

\begin{center}\rule{0.5\linewidth}{0.5pt}\end{center}

\subsection{15.1 Derivata della Softmax rispetto agli
score}\label{derivata-della-softmax-rispetto-agli-score}

Per prima cosa deriviamo \(\hat P^{(l)}\) rispetto a \(s_k\).

\subsubsection{\texorpdfstring{Caso \(l=k\)}{Caso l=k}}\label{caso-lk}

\[
\frac{\partial \hat P^{(k)}}{\partial s_k}
=
\hat P^{(k)}(1-\hat P^{(k)}).
\]

\subsubsection{\texorpdfstring{Caso
\(l\neq k\)}{Caso l\textbackslash neq k}}\label{caso-lneq-k}

\[
\frac{\partial \hat P^{(l)}}{\partial s_k}
=
-\hat P^{(l)}\hat P^{(k)}.
\]

\begin{center}\rule{0.5\linewidth}{0.5pt}\end{center}

\subsection{\texorpdfstring{15.2 Gradiente rispetto a
\(\theta^{(k)}\)}{15.2 Gradiente rispetto a \textbackslash theta\^{}\{(k)\}}}\label{gradiente-rispetto-a-thetak}

Applicando la chain rule, si ottiene il risultato fondamentale:

\[
\nabla_{\theta^{(k)}} L(\theta)
=
\sum_{i=1}^n \big(y_i^{(k)}-\hat P^{(k)}(x_i)\big)\phi(x_i).
\]

Questa formula è l'analogo multiclasse del gradiente logistico binario.

\subsubsection{Messaggio importante}\label{messaggio-importante}

Anche qui la struttura è sempre la stessa:

\[
\text{gradiente} = \text{errore} \times \text{input}.
\]

\begin{center}\rule{0.5\linewidth}{0.5pt}\end{center}

\section{16. Concavità della log-likelihood per
Softmax}\label{concavituxe0-della-log-likelihood-per-softmax}

La dimostrazione completa è più tecnica, ma l'idea è questa: la Hessiana
a blocchi della log-likelihood è semidefinita negativa.

Il modo elegante di vederlo è osservare che, per un singolo esempio, il
termine quadratico può essere riscritto come l'opposto di una varianza,
e quindi è sempre \(\le 0\).

\subsubsection{Conclusione finale}\label{conclusione-finale}

La log-likelihood Softmax è concava. Equivalentemente, la cross-entropy
multiclasse è convessa.

\begin{center}\rule{0.5\linewidth}{0.5pt}\end{center}

\section{17. Perceptron: teorema di
convergenza}\label{perceptron-teorema-di-convergenza}

Assumiamo dati linearmente separabili. Esiste cioè un vettore
\(\theta^*\) con \(\|\theta^*\|=1\) tale che

\[
y_i\,\theta^{*T}x_i \ge \gamma^* > 0
\quad\forall i.
\]

Qui \(\gamma^*\) è il margine di separazione.

L'algoritmo del perceptron aggiorna:

\[
\theta_{t+1}=\theta_t+y_i x_i
\]

quando sbaglia sul punto \((x_i,y_i)\).

Vogliamo mostrare che il numero di errori è finito.

\begin{center}\rule{0.5\linewidth}{0.5pt}\end{center}

\subsection{17.1 Upper bound sulla norma}\label{upper-bound-sulla-norma}

Se facciamo un aggiornamento,

\[
\|\theta_{t+1}\|^2
=
\|\theta_t + y_i x_i\|^2.
\]

Sviluppiamo:

\[
\|\theta_{t+1}\|^2
=
\|\theta_t\|^2 + 2y_i x_i^T\theta_t + \|x_i\|^2.
\]

Nel caso classico del perceptron, l'aggiornamento avviene quando
\(y_ix_i^T\theta_t\le 0\). Quindi

\[
2y_i x_i^T\theta_t \le 0.
\]

Se \(R=\max_i \|x_i\|\), allora

\[
\|\theta_{t+1}\|^2 \le \|\theta_t\|^2 + R^2.
\]

Dopo \(M\) errori,

\[
\|\theta_M\|^2 \le M R^2,
\qquad
\|\theta_M\|\le \sqrt{M}\,R.
\]

\begin{center}\rule{0.5\linewidth}{0.5pt}\end{center}

\subsection{\texorpdfstring{17.2 Lower bound sul prodotto con
\(\theta^*\)}{17.2 Lower bound sul prodotto con \textbackslash theta\^{}*}}\label{lower-bound-sul-prodotto-con-theta}

Consideriamo

\[
\theta^{*T}\theta_{t+1}
=
\theta^{*T}(\theta_t+y_ix_i)
=
\theta^{*T}\theta_t + y_i\theta^{*T}x_i.
\]

Per definizione di margine,

\[
y_i\theta^{*T}x_i \ge \gamma^*.
\]

Quindi ogni aggiornamento aumenta il prodotto di almeno \(\gamma^*\).
Dopo \(M\) errori,

\[
\theta^{*T}\theta_M \ge M\gamma^*.
\]

\begin{center}\rule{0.5\linewidth}{0.5pt}\end{center}

\subsection{17.3 Sandwich finale con
Cauchy-Schwarz}\label{sandwich-finale-con-cauchy-schwarz}

Per Cauchy-Schwarz,

\[
\theta^{*T}\theta_M \le \|\theta^*\|\|\theta_M\| = \|\theta_M\|,
\]

perché \(\|\theta^*\|=1\).

Usando i due bound ottenuti:

\[
M\gamma^* \le \|\theta_M\| \le \sqrt{M}\,R.
\]

Eleviamo al quadrato:

\[
M^2(\gamma^*)^2 \le M R^2.
\]

Se \(M>0\), dividiamo per \(M\):

\[
M \le \frac{R^2}{(\gamma^*)^2}.
\]

\subsubsection{Conclusione}\label{conclusione-8}

Il numero di errori è finito: il perceptron converge in un numero di
passi limitato.

\begin{center}\rule{0.5\linewidth}{0.5pt}\end{center}

\section{18. SVM: dalla formulazione geometrica alla hinge
loss}\label{svm-dalla-formulazione-geometrica-alla-hinge-loss}

\begin{center}\rule{0.5\linewidth}{0.5pt}\end{center}

\subsection{18.1 Hard margin}\label{hard-margin}

Se i dati sono separabili, vogliamo il separatore col margine massimo:

\[
\min_\theta \frac{1}{2}\|\theta\|^2
\quad\text{soggetto a}\quad
y_i\theta^Tx_i \ge 1 \ \forall i.
\]

Minimizzare \(\|\theta\|^2\) equivale a massimizzare il margine
geometrico.

\begin{center}\rule{0.5\linewidth}{0.5pt}\end{center}

\subsection{18.2 Soft margin con variabili
slack}\label{soft-margin-con-variabili-slack}

Se i dati non sono perfettamente separabili, introduciamo slack
\(\xi_i\ge 0\):

\[
\min_{\theta,\xi}
\frac{1}{2}\|\theta\|^2 + C\sum_{i=1}^n \xi_i
\]

soggetto a

\[
y_i\theta^Tx_i \ge 1-\xi_i,
\qquad
\xi_i\ge 0.
\]

\begin{center}\rule{0.5\linewidth}{0.5pt}\end{center}

\subsection{18.3 Come nasce la hinge
loss}\label{come-nasce-la-hinge-loss}

Dal vincolo

\[
y_i\theta^Tx_i \ge 1-\xi_i
\]

otteniamo

\[
\xi_i \ge 1-y_i\theta^Tx_i.
\]

Insieme a \(\xi_i\ge 0\), questo significa che \(\xi_i\) deve essere
almeno il massimo tra queste due quantità:

\[
\xi_i = \max\big(0,\,1-y_i\theta^Tx_i\big).
\]

Questa è esattamente la \textbf{hinge loss}:

\[
(1-y_i\theta^Tx_i)_+.
\]

Sostituendo, otteniamo la formulazione unconstrained:

\[
J(\theta)=\frac{1}{2}\|\theta\|^2 + C\sum_{i=1}^n (1-y_i\theta^Tx_i)_+.
\]

\begin{center}\rule{0.5\linewidth}{0.5pt}\end{center}

\section{19. Sub-gradiente della SVM}\label{sub-gradiente-della-svm}

La hinge loss non è derivabile in \(1-y_i\theta^Tx_i=0\), ma è convessa.
Possiamo quindi usare il \textbf{sub-gradiente}.

\begin{center}\rule{0.5\linewidth}{0.5pt}\end{center}

\subsection{19.1 Caso 1: punto ben classificato con margine
sufficiente}\label{caso-1-punto-ben-classificato-con-margine-sufficiente}

Se

\[
y_i\theta^Tx_i > 1,
\]

allora la hinge vale \(0\) e il contributo del punto alla loss è piatto.
Il sub-gradiente è solo quello del termine di regolarizzazione:

\[
\partial J(\theta)=\theta.
\]

\begin{center}\rule{0.5\linewidth}{0.5pt}\end{center}

\subsection{19.2 Caso 2: punto dentro il margine o classificato
male}\label{caso-2-punto-dentro-il-margine-o-classificato-male}

Se

\[
y_i\theta^Tx_i < 1,
\]

allora

\[
(1-y_i\theta^Tx_i)_+ = 1-y_i\theta^Tx_i.
\]

Derivando rispetto a \(\theta\),

\[
\nabla_\theta(1-y_i\theta^Tx_i) = -y_i x_i.
\]

Quindi il sub-gradiente totale è

\[
\partial J(\theta)=\theta - C y_i x_i.
\]

\begin{center}\rule{0.5\linewidth}{0.5pt}\end{center}

\subsection{19.3 Aggiornamento
stocastico}\label{aggiornamento-stocastico}

Con learning rate \(\gamma\), otteniamo:

\begin{itemize}
\item
  se \(y_i\theta^Tx_i>1\), \[
  \theta \leftarrow \theta - \gamma \theta;
  \]
\item
  se \(y_i\theta^Tx_i<1\), \[
  \theta \leftarrow \theta - \gamma(\theta - C y_i x_i).
  \]
\end{itemize}

Cioè

\[
\theta \leftarrow (1-\gamma)\theta + \gamma C y_i x_i.
\]

\subsubsection{Interpretazione}\label{interpretazione-6}

\begin{itemize}
\tightlist
\item
  il termine \((1-\gamma)\theta\) ``ritira'' i pesi verso zero;
\item
  il termine \(y_i x_i\) corregge l'errore di classificazione.
\end{itemize}

\begin{center}\rule{0.5\linewidth}{0.5pt}\end{center}

\section{20. Kernel trick e argomento di
rappresentazione}\label{kernel-trick-e-argomento-di-rappresentazione}

Supponiamo di mappare i dati in uno spazio di feature alto o infinito
tramite \(\Psi(x)\).

Nel nuovo spazio, il classificatore lineare è

\[
f(x)=\theta^T\Psi(x).
\]

L'idea chiave è che la soluzione ottima \(\theta\) vive nello span dei
dati trasformati.

\begin{center}\rule{0.5\linewidth}{0.5pt}\end{center}

\subsection{\texorpdfstring{20.1 Perché \(\theta\) può essere scritto
come combinazione dei
dati}{20.1 Perché \textbackslash theta può essere scritto come combinazione dei dati}}\label{perchuxe9-theta-puuxf2-essere-scritto-come-combinazione-dei-dati}

Scomponiamo \(\theta\) in due parti:

\[
\theta = \theta_{\parallel} + \theta_{\perp},
\]

dove:

\begin{itemize}
\tightlist
\item
  \(\theta_{\parallel}\) sta nello spazio generato dai vettori
  \(\Psi(x_1),\dots,\Psi(x_n)\);
\item
  \(\theta_{\perp}\) è ortogonale a quello spazio.
\end{itemize}

Per ogni training point \(x_i\),

\[
\theta^T\Psi(x_i)
=
\theta_{\parallel}^T\Psi(x_i) + \theta_{\perp}^T\Psi(x_i).
\]

Ma \(\theta_{\perp}\) è ortogonale a tutto lo span, quindi

\[
\theta_{\perp}^T\Psi(x_i)=0.
\]

Dunque la classificazione dipende solo da \(\theta_{\parallel}\).

Invece la norma si scompone come

\[
\|\theta\|^2 = \|\theta_{\parallel}\|^2 + \|\theta_{\perp}\|^2.
\]

Per minimizzare la norma, conviene mettere \(\theta_{\perp}=0\).

\subsubsection{Conclusione}\label{conclusione-9}

La soluzione ottima ha la forma

\[
\theta = \sum_{i=1}^n \alpha_i \Psi(x_i).
\]

\begin{center}\rule{0.5\linewidth}{0.5pt}\end{center}

\subsection{20.2 Il kernel}\label{il-kernel}

Sostituendo nella funzione di decisione,

\[
f(x)
=
\left(\sum_{i=1}^n \alpha_i \Psi(x_i)\right)^T \Psi(x)
=
\sum_{i=1}^n \alpha_i \langle \Psi(x_i),\Psi(x)\rangle.
\]

Definiamo quindi il kernel

\[
K(x_i,x)=\langle \Psi(x_i),\Psi(x)\rangle.
\]

Otteniamo

\[
f(x)=\sum_{i=1}^n \alpha_i K(x_i,x).
\]

\subsubsection{Interpretazione}\label{interpretazione-7}

Non serve conoscere esplicitamente \(\Psi(x)\): basta saper calcolare
\(K(x,z)\).

\begin{center}\rule{0.5\linewidth}{0.5pt}\end{center}

\section{21. Statistical Learning Theory: classi finite e sample
complexity}\label{statistical-learning-theory-classi-finite-e-sample-complexity}

Sia \(H\) una classe finita di classificatori. Vogliamo sapere quanti
campioni servono per garantire che l'errore empirico sia vicino
all'errore vero \textbf{uniformemente} su tutti gli \(h\in H\).

\begin{center}\rule{0.5\linewidth}{0.5pt}\end{center}

\subsection{21.1 Hoeffding per un classificatore
fissato}\label{hoeffding-per-un-classificatore-fissato}

Per un classificatore fissato \(h\), l'errore empirico \(\hat R(h)\) è
la media di variabili Bernoulli indipendenti. Per Hoeffding,

\[
P\big(|\hat R(h)-R(h)|>\varepsilon\big)
\le
2e^{-2n\varepsilon^2}.
\]

\begin{center}\rule{0.5\linewidth}{0.5pt}\end{center}

\subsection{21.2 Union bound su tutta la
classe}\label{union-bound-su-tutta-la-classe}

A noi non basta controllare un solo \(h\): vogliamo che il bound valga
per \textbf{tutti} gli \(h\in H\).

Usiamo quindi l'union bound:

\[
P\left(\sup_{h\in H}|\hat R(h)-R(h)|>\varepsilon\right)
\le
\sum_{h\in H} P\big(|\hat R(h)-R(h)|>\varepsilon\big).
\]

Ogni termine è limitato da \(2e^{-2n\varepsilon^2}\), quindi

\[
P\left(\sup_{h\in H}|\hat R(h)-R(h)|>\varepsilon\right)
\le
2|H|e^{-2n\varepsilon^2}.
\]

\begin{center}\rule{0.5\linewidth}{0.5pt}\end{center}

\subsection{21.3 Ricavare la sample
complexity}\label{ricavare-la-sample-complexity}

Se vogliamo che questa probabilità sia al più \(\delta\), imponiamo

\[
2|H|e^{-2n\varepsilon^2}\le \delta.
\]

Prendiamo il logaritmo:

\[
\log 2 + \log |H| - 2n\varepsilon^2 \le \log \delta.
\]

Equivalentemente,

\[
2n\varepsilon^2 \ge \log \frac{2|H|}{\delta}.
\]

Quindi basta prendere

\[
n \ge \frac{1}{2\varepsilon^2}\log\frac{2|H|}{\delta}.
\]

\subsubsection{Conclusione}\label{conclusione-10}

La sample complexity cresce solo \textbf{logaritmicamente} con la
cardinalità della classe \(H\).

\begin{center}\rule{0.5\linewidth}{0.5pt}\end{center}

\section{22. VC dimension: esempi
classici}\label{vc-dimension-esempi-classici}

La VC dimension di una classe \(H\) è il massimo numero di punti che
\(H\) può \textbf{shatterare}, cioè etichettare in tutti i modi
possibili.

\begin{center}\rule{0.5\linewidth}{0.5pt}\end{center}

\subsection{22.1 Intervalli sulla retta: VC dimension =
2}\label{intervalli-sulla-retta-vc-dimension-2}

\subsubsection{Perché almeno 2}\label{perchuxe9-almeno-2}

Prendiamo due punti \(x_1<x_2\). Un intervallo \([a,b]\) può realizzare
tutte le etichettature:

\begin{itemize}
\tightlist
\item
  \(--\): scegli intervallo vuoto;
\item
  \(+-\): prendi un intervallo che contiene solo \(x_1\);
\item
  \(-+\): prendi un intervallo che contiene solo \(x_2\);
\item
  \(++\): prendi un intervallo che contiene entrambi.
\end{itemize}

Quindi 2 punti sono shatterabili.

\subsubsection{Perché non 3}\label{perchuxe9-non-3}

Prendiamo tre punti ordinati \(x_1<x_2<x_3\). Considera l'etichettatura

\[
+,-,+.
\]

Nessun intervallo singolo può includere \(x_1\) e \(x_3\) senza
includere anche \(x_2\).

Quindi 3 punti non sono shatterabili.

\subsubsection{Conclusione}\label{conclusione-11}

\[
d_{VC}=2.
\]

\begin{center}\rule{0.5\linewidth}{0.5pt}\end{center}

\subsection{\texorpdfstring{22.2 Iperpiani in \(\mathbb{R}^2\): VC
dimension =
3}{22.2 Iperpiani in \textbackslash mathbb\{R\}\^{}2: VC dimension = 3}}\label{iperpiani-in-mathbbr2-vc-dimension-3}

\subsubsection{Perché almeno 3}\label{perchuxe9-almeno-3}

Tre punti in posizione generale nel piano possono essere separati in
tutti i modi da una retta.

\subsubsection{Perché non 4}\label{perchuxe9-non-4}

Prendiamo quattro punti in configurazione XOR. L'etichettatura alternata
non è linearmente separabile.

\subsubsection{Conclusione}\label{conclusione-12}

Le rette in \(\mathbb{R}^2\) hanno

\[
d_{VC}=3.
\]

Più in generale, gli iperpiani in \(\mathbb{R}^d\) hanno VC dimension
\(d+1\).

\begin{center}\rule{0.5\linewidth}{0.5pt}\end{center}

\section{23. Decision Trees: perché uno split non peggiora
l'errore}\label{decision-trees-perchuxe9-uno-split-non-peggiora-lerrore}

Consideriamo una foglia \(l\) che contiene \(N_l\) punti, di cui
\(N_l^+\) positivi.

Definiamo

\[
z_l = \frac{N_l^+}{N_l}.
\]

L'impurità del nodo è una funzione \(\psi(z_l)\), ad esempio:

\begin{itemize}
\tightlist
\item
  errore di classificazione;
\item
  indice di Gini;
\item
  entropia.
\end{itemize}

L'errore totale del nodo è

\[
E(l)=N_l\psi(z_l).
\]

Supponiamo di dividere il nodo in figlio sinistro \(L\) e figlio destro
\(R\).

\begin{center}\rule{0.5\linewidth}{0.5pt}\end{center}

\subsection{23.1 La media della madre è combinazione convessa delle
medie dei
figli}\label{la-media-della-madre-uxe8-combinazione-convessa-delle-medie-dei-figli}

Poiché

\[
N_l=N_L+N_R,
\qquad
N_l^+=N_L^+ + N_R^+,
\]

abbiamo

\[
\frac{N_l^+}{N_l}
=
\frac{N_L^+ + N_R^+}{N_l}.
\]

Moltiplichiamo e dividiamo nei termini giusti:

\[
\frac{N_l^+}{N_l}
=
\frac{N_L}{N_l}\frac{N_L^+}{N_L}
+
\frac{N_R}{N_l}\frac{N_R^+}{N_R}.
\]

Se definiamo

\[
\alpha=\frac{N_L}{N_l},
\qquad
1-\alpha=\frac{N_R}{N_l},
\]

allora

\[
z_l = \alpha z_L + (1-\alpha)z_R.
\]

Quindi l'impurità della madre è valutata in una combinazione convessa
delle impurità ``grezze'' dei figli.

\begin{center}\rule{0.5\linewidth}{0.5pt}\end{center}

\subsection{23.2 Uso della concavità}\label{uso-della-concavituxe0}

Le funzioni di impurità standard sono \textbf{concave} in \(z\). Per la
concavità,

\[
\psi(\alpha z_L + (1-\alpha)z_R)
\ge
\alpha\psi(z_L)+(1-\alpha)\psi(z_R).
\]

Sostituendo le definizioni,

\[
\psi\left(\frac{N_l^+}{N_l}\right)
\ge
\frac{N_L}{N_l}\psi\left(\frac{N_L^+}{N_L}\right)
+
\frac{N_R}{N_l}\psi\left(\frac{N_R^+}{N_R}\right).
\]

Moltiplichiamo per \(N_l\):

\[
N_l\psi\left(\frac{N_l^+}{N_l}\right)
\ge
N_L\psi\left(\frac{N_L^+}{N_L}\right)
+
N_R\psi\left(\frac{N_R^+}{N_R}\right).
\]

Cioè

\[
E(\text{madre}) \ge E(\text{figlio sinistro}) + E(\text{figlio destro}).
\]

\subsubsection{Conclusione}\label{conclusione-13}

Uno split non aumenta mai l'errore impurità totale.

\begin{center}\rule{0.5\linewidth}{0.5pt}\end{center}

\section{24. Random Forests: perché
aiutano}\label{random-forests-perchuxe9-aiutano}

Un singolo albero è ad alta varianza: piccole variazioni nei dati
possono cambiare molto la struttura appresa.

Le Random Forests riducono questa instabilità combinando:

\begin{enumerate}
\def\labelenumi{\arabic{enumi}.}
\tightlist
\item
  \textbf{bootstrap} sui campioni;
\item
  \textbf{randomizzazione delle feature} ad ogni split;
\item
  \textbf{media / voto} di molti alberi.
\end{enumerate}

\subsubsection{Effetto bias-varianza}\label{effetto-bias-varianza}

\begin{itemize}
\tightlist
\item
  il bias non aumenta troppo;
\item
  la varianza si riduce parecchio.
\end{itemize}

Per questo le Random Forests generalizzano meglio di un singolo albero
profondo.

\begin{center}\rule{0.5\linewidth}{0.5pt}\end{center}

\section{25. Riassunto finale super-compatto per
l'esame}\label{riassunto-finale-super-compatto-per-lesame}

\subsection{Regressione e predittori
ottimali}\label{regressione-e-predittori-ottimali}

\begin{itemize}
\tightlist
\item
  loss quadratica \(\Rightarrow\) media condizionata;
\item
  loss assoluta \(\Rightarrow\) mediana condizionata;
\item
  classificazione binaria 0-1 \(\Rightarrow\) classe a posteriori più
  probabile.
\end{itemize}

\subsection{OLS}\label{ols}

\begin{itemize}
\tightlist
\item
  minimizza \(\|y-X\theta\|^2\);
\item
  soluzione: \[
  \hat\theta=(X^TX)^{-1}X^Ty;
  \]
\item
  residui ortogonali: \[
  X^T(y-X\hat\theta)=0;
  \]
\item
  Hessiana: \[
  X^TX \succeq 0;
  \]
\item
  OLS è unbiased: \[
  \mathbb{E}[\hat\theta]=\theta^*;
  \]
\item
  varianza: \[
  \operatorname{Var}(\hat\theta)=\sigma^2(X^TX)^{-1}.
  \]
\end{itemize}

\subsection{Ridge}\label{ridge}

\begin{itemize}
\tightlist
\item
  aggiunge \(\lambda\|\theta\|^2\);
\item
  riduce varianza;
\item
  soluzione: \[
  \hat\theta_{\text{ridge}}=(X^TX+\lambda I)^{-1}X^Ty.
  \]
\end{itemize}

\subsection{Logistic Regression}\label{logistic-regression}

\begin{itemize}
\tightlist
\item
  nasce da Bernoulli + MLE;
\item
  likelihood: \[
  \prod_i p_i^{y_i}(1-p_i)^{1-y_i};
  \]
\item
  objective equivalenti:

  \begin{itemize}
  \tightlist
  \item
    massimizzare log-likelihood;
  \item
    minimizzare negative log-likelihood;
  \item
    minimizzare cross-entropy.
  \end{itemize}
\item
  Hessiana negativa \(\Rightarrow\) log-likelihood concavo.
\end{itemize}

\subsection{Softmax}\label{softmax}

\begin{itemize}
\tightlist
\item
  gradiente: \[
  \nabla_{\theta^{(k)}}L
  =
  \sum_i (y_i^{(k)}-\hat P^{(k)}(x_i))\phi(x_i).
  \]
\end{itemize}

\subsection{Perceptron}\label{perceptron}

\begin{itemize}
\tightlist
\item
  converge se i dati sono linearmente separabili;
\item
  numero di errori limitato da \(R^2/(\gamma^*)^2\).
\end{itemize}

\subsection{SVM}\label{svm}

\begin{itemize}
\tightlist
\item
  massimizza il margine;
\item
  soft margin \(\Rightarrow\) hinge loss;
\item
  kernel trick \(\Rightarrow\) lavora con prodotti scalari in spazi di
  feature impliciti.
\end{itemize}

\subsection{Statistical Learning
Theory}\label{statistical-learning-theory}

\begin{itemize}
\tightlist
\item
  classi finite: \[
  n \ge \frac{1}{2\varepsilon^2}\log\frac{2|H|}{\delta};
  \]
\item
  la VC dimension misura la complessità della classe.
\end{itemize}

\subsection{Decision Trees}\label{decision-trees}

\begin{itemize}
\tightlist
\item
  lo split non peggiora l'impurità grazie alla concavità di Gini /
  Entropia;
\item
  le Random Forests riducono la varianza del singolo albero.
\end{itemize}

\begin{center}\rule{0.5\linewidth}{0.5pt}\end{center}

\section{26. Osservazioni finali sulle confusioni più
comuni}\label{osservazioni-finali-sulle-confusioni-piuxf9-comuni}

\subsection{1. ``Linear regression'' non significa ``retta nel
grafico''}\label{linear-regression-non-significa-retta-nel-grafico}

Significa \textbf{lineare nei parametri}. Posso usare feature come
\(x, x^2, x^3\) e ottenere curve non lineari in \(x\), restando lineare
in \(\theta\).

\subsection{2. Logistic regression ha ``due
obiettivi''?}\label{logistic-regression-ha-due-obiettivi}

No: di solito è lo stesso obiettivo scritto in modi diversi.

\[
\max \ell(\theta)
\quad\Longleftrightarrow\quad
\min -\ell(\theta).
\]

\subsection{\texorpdfstring{3. Perché compare \(X^T\) in
OLS?}{3. Perché compare X\^{}T in OLS?}}\label{perchuxe9-compare-xt-in-ols}

Perché il residuo deve essere ortogonale allo spazio delle colonne di
\(X\), e questo si scrive

\[
X^T(y-X\hat\theta)=0.
\]

\subsection{4. Boosting riduce il bias o la
varianza?}\label{boosting-riduce-il-bias-o-la-varianza}

Principalmente \textbf{riduce il bias}, perché aggiunge modelli deboli
in sequenza per correggere errori sistematici dei precedenti.

\subsection{5. Bagging e Random Forest cosa
riducono?}\label{bagging-e-random-forest-cosa-riducono}

Principalmente \textbf{la varianza}, mediando molti modelli instabili.

\begin{center}\rule{0.5\linewidth}{0.5pt}\end{center}

%%%%%%%%%%%%%%%%%%%%%%%%%%%%%%%%%%%%%%%%%%%%%%%%%%%%%%%%%%%%


\end{document}
